\section{Conclusion}
\label{sec:conclusion}
This paper characterises data-centre flexibility through an integrated representation of IT workload, UPS storage and cooling with TES. The resulting response depends on the requested power deviation, its start time and the operating state inherited by each resource. Including recovery connects event delivery to the facility's subsequent obligations, while component decomposition reveals the internal actions behind the metered response.

For the 1 MW IT-rated case, the representative-day assessment demonstrates a 100 kW import reduction lasting six hours and a 25 kW import increase reaching the twelve-hour assessment cap. The participating resources vary with timing: selected morning reductions are predominantly supplied by UPS discharge, whereas afternoon reductions rely mainly on workload deferral and reduced cooling. These examples show why equipment capacity alone is insufficient to describe available flexibility, and why a whole-facility assessment is needed to capture compensating actions across assets.

Annual coordinated operation reduces electricity cost by GBP~29,626.91, or 5.38\%, under the observed 2025 reference prices. This saving accompanies greater annual electricity use and peak import, demonstrating the importance of distinguishing price value from energy efficiency and network benefit. Flexible workload share produces the largest economic response among the tested assumptions. Enlarging UPS capacity from 600 to 900 kWh adds approximately GBP~2,460/year in savings and supports a conditional investment threshold, rather than an unconditional case for larger batteries.

The findings support describing flexibility as a state-dependent capability with magnitude, duration and recovery conditions. Extending the assessment to varied workloads and weather, explicit degradation and tariffs, and a broader set of event days would support site-specific investment and market-participation decisions. The present results provide the physical and operational basis for those decisions while keeping annual electricity-cost savings separate from unmodelled service revenues.
